% Ad Familiares 3.11 (ad-familiares-03-11) --- English translation
% Translated by Claude (Anthropic) from the Latin (Perseus).
% Style guide: STYLE.md.

\ciceroLetterOpener{CICERO APPIO PVLCHRO VT SPERO, CENSORI S. D.}

\ciceroSection{1} When I was in camp by the river Pyramus, two letters from you were brought to me at the same time, which Q. Servilius had sent me from Tarsus. The one had the date marked on it: the Nones of April; the other, which seemed to me the more recent, was undated. I will answer the earlier first, in which you write to me of your acquittal on the \textit{maiestas} charge. Of this, although I had been informed long before --- by letters, by messengers, by rumour itself in the end (for nothing was more public; not that anyone had thought the other way, but where men distinguished for their renown are concerned nothing is reported in muffled terms) --- still, that same news was made the more cheerful for me by your letter, not only because it spoke more plainly and more fully than common talk, but also because I felt I was congratulating you the more truly when I heard about you from yourself.

\ciceroSection{2} So I embraced you in thought even in your absence; the letter itself I actually kissed, and I congratulated myself as well. For what the whole people, the Senate, and the jurors render in tribute to your talent, your industry, your virtue --- because perhaps I flatter myself, when I imagine those qualities to be in me --- I count as rendered to me too. Nor did I marvel that the outcome of your trial was so glorious, but that your enemies' minds were so warped. ``But what difference does it make,'' you will say, ``\textit{ambitus} or \textit{maiestas}?'' For the substance, none: the one you have not touched, the other you have positively magnified. \textit{Maiestas} indeed, however, is what Sulla meant it to be, so that men should not be free to declaim against anyone they pleased with impunity; but \textit{ambitus} carries with it so plain a violence that the charge must either be brought disgracefully or be defended. For: can it really be unknown whether or not bribery has been committed? And in your case, whose suspicions has the course of your honours ever raised?

\ciceroSection{3} Wretched me, that I was not there! What peals of laughter I would have raised! But on the \textit{maiestas} verdict, two things in your letter gave me the greatest pleasure: one, that you write you were defended by the commonwealth itself --- which even amid an abundance of good and brave citizens ought to be standing by men of your stamp, and far more now, when there is such a dearth at every rank, of office or of age, that a state thus orphaned ought to embrace its few remaining guardians; the other, that you praise the loyalty and goodwill of Pompey and of Brutus in marvellous terms. I rejoice at the virtue and the dutifulness both of your kinsmen, who are also my closest friends --- the one of them the first man of every age and nation, the other long since the first of our young men and soon, as I hope, of our state. As for marking the hired witnesses by way of their own communities: unless something has already been done in the matter through Flaccus, I shall see to it as I make my way back through Asia.

\ciceroSection{4} Now I come to the second letter. That you have sent me, as it were, a likeness of the common condition of the times and of the whole commonwealth, drawn out in detail, is most welcome to me on account of the practical wisdom of your letter; for I see both that the dangers are lighter than I feared, and that the resources are greater --- if, as you write, all the forces of the state have ranged themselves under Pompey's command. At the same time I have read clear evidence of your own prompt and eager spirit in the defence of the commonwealth, and I have taken extraordinary pleasure from this attentiveness of yours --- that in the midst of your highest occupations you should yet have wished the state of the commonwealth to be known to me through you. As for the books on augural law, keep them for the leisure we shall share. For when by letter I was pressing you to deliver on your promises, I supposed you to be most at leisure near the city; now, however, as you yourself pledge, in place of the augural books I shall await all your speeches collected.

\ciceroSection{5} D. Tullius, to whom you gave commissions for me, had not met up with me; and there was no one of your people now with me apart from all my own --- who are all yours. As to which of my letters you call ``rather peevish,'' I do not understand. Twice I wrote to you, carefully exonerating myself, and faulting you mildly for having believed, on no evidence, things about me. This kind of complaint did seem to me, indeed, the part of a friend; but if it displeases you, I will not use it hereafter. As for that letter not being ``eloquent,'' as you put it: if so, you may be sure it was not mine. For as Aristarchus pronounces any line he does not approve to be not Homer's, so you in the same way --- for I like to joke --- if anything is unpolished, do not suppose it mine. Farewell, and in your censorship, if you are now censor, as I hope, think much of your great-grandfather.
